%------------------------------------------------------------------------------%
\section{Avaliando Formas Indeterminadas}

Outra operação que pode ser complexa é o cálculo de limites indeterminados,
nesse caso a Série de Taylor também pode ajudar.
Vamos ilustrar a técnica caculando o limite
\[
\lim_{x\to 1} \frac{\ln x}{\;x-1\;} \;.
\]
Esse é um limite, com indeterminação do tipo $\sfrac{0}{0}$,
onde podemos utilizar a regra de l'Hôpital
\[
\lim_{x\to 1} \frac{\ln x}{\;x-1\;}
= \lim_{x\to 1} \frac{\;\sfrac{1}{x}\;}{1}
= 1 \;.
\]
Entretanto, a Série de Taylor oferece uma forma alternativa
para calcular esse limite que pode ser útil em outras situações.

\begin{example}
  Calcule o limite
  \(
  \lim_{x\to 1} \frac{\ln x}{\;x-1\;}
  \).
  \solution
  Primeiro obtemos a Série de Taylor da função $\ln(x)$ com centro em $a=1$
  \[
  \ln(x) = (x-1) - \frac{(x-1)^2}{2}
  + \frac{(x-1)^3}{3}
  - \frac{(x-1)^4}{4}
  + \cdots
  \]
  a igualdade é verdadeira para todo $x$ tal que $\abs{x-1} < 1$.
  Agora substituímos o logaritmo pela sua série na função original
  \[
  \frac{\ln x}{\,x-1\,}
  = \frac{1}{\,x-1\,}
  \left( (x-1)
  - \frac{(x-1)^2}{2}
  + \frac{(x-1)^3}{3}
  - \frac{(x-1)^4}{4}
  + \cdots
  \right) \;.
  \]
  Manipulamos os termos da série obtemos
  \[
  \frac{\ln x}{\,x-1\,}
  = 1 - \frac{x-1}{2} + \frac{(x-1)^2}{3} - \frac{(x-1)^3}{4} + \cdots
  \]
  Agora calculamos o limite desejado
  \[
  \lim_{x\to 1} \frac{\ln x}{\,x-1\,}
  = \lim_{x\to 1} \left( 1
  - \frac{x-1}{2}
  + \frac{(x-1)^2}{3}
  - \frac{(x-1)^3}{4}
  + \cdots
  \right)
  \]
  como todos os termos da série, com exceção do primeiro, vão para zero,
  quando $x\to 1$, temos que
  \[
  \lim_{x\to 1} \frac{\ln x}{\,x-1\,} = 1 \;.
  \]
\end{example}

%------------------------------------------------------------------------------%
